\chapter{Reproducibility When the Tool Varies}
\label{ch:reproducibility}

Reproducibility covers several distinct questions, described in
\Cref{tab:reproducibility}. The words are not used the same way everywhere.
\textcite{barba2018terminologies} found two camps that swap
\emph{reproduce} and \emph{replicate}, and recommends defining the terms you
use. The table follows the more common camp: reproducing uses the same data
and code, replicating uses new data. Language models add a difficulty to these
familiar questions: a step in the pipeline may give a different answer each
time it runs.

\begin{table}[htbp]
\centering
\caption{Terms of art, and what each one assesses. The last row is not a
standard term; it is the version of the reference data a result was computed
against.}
\label{tab:reproducibility}
\footnotesize
\begin{tabularx}{\textwidth}{lXXXX}
\toprule
\textbf{Term} & \textbf{Definition} & \textbf{Held constant} & \textbf{Varied} & \textbf{Assesses} \\
\midrule
Repeatability & Same team, same result, identical conditions & Team, data, code, environment & Nothing & Internal consistency \\
\addlinespace
Reproducibility & Independent team, same result from same data and code & Data, code & Team, environment & Computational integrity \\
\addlinespace
Replicability & Independent teams, comparable conclusions from new data & Question, method & Team, infrastructure, data & Scientific validity \\
\addlinespace
Reusability & Data or code repurposed for other ends & Digital assets & Purpose, context & Long-term usefulness \\
\addlinespace
Provenance & Documented origin and processing history & Metadata, logs & Execution snapshot & Transparency \\
\addlinespace
Reference version & The backbone, reference database or annotation release used & Data, code, environment & Release of the reference & Dependence on the reference \\
\bottomrule
\end{tabularx}
\end{table}

\section{Where a language model breaks this}

\textbf{Sampling is random.} The same input can give different outputs. Setting
the temperature to zero reduces the variation but does not remove it
\parencite{ouyang2024nondeterminism,coqueret2026randomness}. What remains
comes from how a server batches requests, from the hardware and from silent
model updates. Several hosted reasoning models do not accept a temperature at
all, and always taking the most probable token can cost performance. A seed
repeats the random draws only while everything else stays the same
\parencite{coqueret2026randomness}. I use temperature zero where a model
allows it, unless its developers advise against it. Whatever the settings,
record them (\Cref{fig:variation}), keep the raw outputs, and rerun a sample
to see how much they vary. \textcite{keck2026interactions} reran 100
paragraphs 100 times each at temperature zero, and the runs did not always
agree.

\textbf{The model changes underneath you.} A hosted model with a fixed name is
not a fixed thing; providers update it without notice. A local model file with a
known checksum does not change --- the strongest argument for open weights in
research. The file is not the whole instrument, though. Outputs also depend on
the program that runs it, the numerical precision, the processor and its
drivers, and how requests are batched, so these go in the log too
\parencite{coqueret2026randomness}.

\textbf{A generated claim has no provenance.} Keep the model out of the chain of
inference. Use it to write code or a query, which you read and run; the code is
the artefact you keep. Do not let it produce a number that goes into a result,
unless you have measured how often it is wrong. That is the third path in
\Cref{fig:chain}: the model as an instrument, whose output enters the analysis
only through a validation step (\Cref{sec:outputasdata}).

\begin{figure}[htbp]
\centering
\fig{ch09-chain-of-inference}
\caption{Three places for a model in an analysis. Generated code is read, run
and committed, so the result can be reproduced. When the model's output is the
data, it enters the analysis only through a sample checked by people, and its
error rates go with it (\Cref{sec:outputasdata}). A generated number that skips
both has no provenance.}
\label{fig:chain}
\end{figure}

\textbf{Agents take different paths}, and may report a task done when it is not
\parencite{shapira2026agentsofchaos}. The deliverable is the reviewed, committed
code, not the session.

\textbf{Context is hidden state.} The same prompt late in a conversation behaves
differently. For anything that matters, use a fresh session with a
self-contained prompt, and save it.

\begin{figure}[htbp]
\centering
\fig{ch09-variation}
\caption{Five sources of variation in a model call, and what to record or
change for each.}
\label{fig:variation}
\end{figure}

\begin{keypoint}
Generated code is reproducible; generated answers are not. Move every
model-dependent step as far upstream as it will go, and keep the deterministic
artefact.
\end{keypoint}

\section{Log model calls like instrument readings}

\begin{lstlisting}[caption={logs/2026-10-15\_taxon-matching.md}]
Model: qwen3-coder:30b, Q4_K_M, sha256:4f2a...c91
Runtime: ollama 0.31.1   Harness: opencode 1.17.11   Context: 64K
Temperature: 0   Seed: 42   Fresh session   Rules: AGENTS.md @ 7c1e2d0
Prompt: Write an R script matching data/raw/species_list.csv against
  the GBIF backbone; keep unmatched rows as status="NO_MATCH".
Reference: GBIF Backbone Taxonomy as of 2026-10-15, DOI in README.md
Output: scripts/03_match_taxa.R (commit a3f9c21)
Review: 1,847 rows in, 1,847 out. 23 NO_MATCH checked by hand:
  19 typos fixed in a lookup table, 4 names not in the backbone
  referred to a specialist.
\end{lstlisting}

Each line records something that can change the result: the runtime and
harness versions, the quantisation, the context length, the rules file the
harness read, and the backbone the names were matched against. For a model
with a reasoning setting, log that too. Keep this record in plain Markdown
alongside the analysis, and commit it with the code. A short record is enough
if it lets someone check what was done.

\section{When the model's output is the data}
\label{sec:outputasdata}

The rule above, that a model should not produce a number that goes into a
result, cannot cover every use. The extraction studies in
\Cref{ch:applications} use a model to read traits, incidences and species
names out of text, and the table it fills is the data. The model is then a
measuring instrument, and an instrument is trusted only after it has been
checked against a known standard. This is the third path in \Cref{fig:chain}.

The check is a validation sample. Draw records at random from everything the
model processed, have a person extract the same fields from the sources, and
compare the two. Where some groups are rare, such as a taxon or a kind of
source, stratify the sample so that each group is checked.
\textcite{keck2026interactions} validated a workflow that extracts species
interactions from the literature in this way. One ecologist annotated 500
paragraphs chosen at random, and the workflow found 70 per cent of the
interactions the annotator found (95 per cent interval 63 to 77). A second,
stratified set covered rare taxa, but the authors note that it was drawn from
the workflow's own output, so it cannot show what the workflow missed.

Report the error of each field separately, with its interval, because a model
can be reliable on one field and wrong on the next. In the plant-disease
reports of \textcite{gougherty2024extraction}, names were mostly right, but
the model wrote 100 per cent incidence for 53 reports that gave none
(\Cref{ch:applications}). Such errors are not random noise; this one pushes
incidence up. Agents that curate metadata make the same kind of error.
Outside ecology, \textcite{youngblut2026scbasecount} had an agent label
single-cell datasets from the Sequence Read Archive. They checked its labels
against an independently curated repository for the datasets the two share.
They also report that a study's disease label can be copied to its healthy
controls.

The error then has to be carried into the analysis. Treating model output as
if it were measured without error gives intervals that are too narrow: in the
Amazon deforestation example of \textcite{angelopoulos2023ppi}, such an
interval missed the true value. Their prediction-powered inference, and the
estimator that \textcite{egami2023surrogates} developed for language-model
labels in social science, use the validated sample to correct estimates made
from all the model's outputs. Egami and colleagues need the chance of each
record entering the validation sample to be known, one more reason to draw it
at random. Where neither method fits the analysis, rerun it on the validated
records alone, or with the measured error rates applied, and report whether
the conclusion changes.

Last, mark the records before they leave the project. Records that reach GBIF
or OBIS without a flag cannot be told apart from records a person made, and
Darwin Core has terms for this \parencite{wieczorek2012darwincore}.
\code{basisOfRecord} still says what the original record was: one read from a
paper is a human observation or a material citation, not a machine
observation, because the model read it rather than made it. Where the model
assigned the name, \code{identifiedBy} can name it with its version, and
\code{identificationVerificationStatus} can say that no person has checked it.
Where the model copied an author's identification, say so in the record's
remarks and in the dataset's methods.

\section{The methods section}

Nature Portfolio's policy, as \textcite{naturemethods2026responsibly}
describe it, asks authors who use AI to help write a manuscript to declare it
in the Methods section and to check the generated text. A model cannot be an
author, because it cannot take responsibility for the content.
\textcite{coqueret2026randomness} ask journals to require more when a model
generates data: the model name and version, its settings, the full prompts,
the date of each query, and the raw outputs in the replication package. Two
templates, one for code and one for data:

\begin{quote}\small
Data processing scripts were drafted with \textit{[model, version]} via
\textit{[tool]} in \textit{[month year]}. All generated code was reviewed,
run locally and validated by \textit{[check]}. Prompts and settings are archived
at \textit{[DOI]}. The authors are fully responsible for the content.

Fields \textit{[names]} were extracted from \textit{[sources]} with
\textit{[model, version]} at \textit{[settings]} in \textit{[month year]}. A
random sample of \textit{[n]} records, stratified by \textit{[groups]}, was
checked against the sources by \textit{[person]}; precision and recall per
field, with 95\% intervals, are in \textit{[table]}. Errors were handled by
\textit{[correction or sensitivity analysis]}. Raw model outputs are archived
at \textit{[DOI]}.
\end{quote}

\section{Containers}

A container pins the operating system, R and its packages. It does not remove
model randomness, and it needs no model at all when the model only wrote code:
the committed scripts are the pipeline. The image below restores the packages
recorded by \code{renv}, and refuses to build if a raw file no longer matches
the checksums written in \Cref{ch:organisation}:

\begin{lstlisting}[style=shell,caption={Dockerfile}]
FROM docker.io/rocker/r-ver:4.4.1
WORKDIR /project
# rocker/r-ver does not include renv
RUN R -q -e 'install.packages("renv")'
COPY renv.lock renv.lock
RUN R -q -e 'renv::restore()'
COPY data/raw.sha256 data/raw.sha256
COPY data/raw data/raw
# the build stops here if a raw file differs from its checksum
RUN sha256sum -c data/raw.sha256
COPY scripts scripts
\end{lstlisting}

Comments go on lines of their own: after an instruction, Docker reads a
\code{\#} and the words after it as more arguments. Pin the image to the R
version recorded in \code{renv.lock}; the tag above is the one tested
(\Cref{app:stateofplay}).

When the model is the instrument (\Cref{sec:outputasdata}), the run also needs
the weights and a program that runs them, such as llama.cpp or Ollama
(\Cref{ch:stack}); this image includes neither. Keep the weights outside the
image, check them against \code{models.sha256}, written as in
\Cref{ch:organisation}, mount them read-only and turn the network off:

\begin{lstlisting}[style=shell]
docker build -t posidonia .
shasum -a 256 -c models.sha256
docker run --rm --network none -v "$PWD/models:/project/models:ro" \
  posidonia Rscript scripts/main.R
\end{lstlisting}

\code{podman} takes the same commands. Where the licence allows, deposit the
weights and their checksum in an archive such as Zenodo, so that the
instrument itself can be cited. Open weights on a local stack you control give
the best chance of exact reproduction, and only if the whole stack is recorded
\parencite{coqueret2026randomness}. A workflow built on a hosted API can be
documented, but not reproduced, because the instrument is not under your
control.
